Recent transformer and foundation-model advances have made large language models (LLMs) increasingly capable of flexible multi-step planning and decision support, which makes them attractive for simulation-centric scenario exploration workflows \citep{vaswani2017attention,devlin2019bert,brown2020llm,openai2023gpt4,touvron2023llama,touvron2023llama2,zhao2023llmsurvey,yang2024chatgptsurvey}. In the military simulation domain specifically, recent work has explored LLM-based commander agents for warfare simulation and tactical wargaming \citep{liu2026commandagent,geo2026commander,liu2026tacticalwargame}, while comprehensive benchmarks have been proposed to evaluate LLM readiness for military decision-making \citep{warbench2026}. Reasoning-time performance can be further shaped by prompt-side strategies such as chain-of-thought, self-consistency, ReAct-style tool interaction, deliberate search over thought branches, and explicit tool use \citep{wei2022cot,wang2023selfconsistency,yao2023react,yao2023tree,schick2023toolformer}. At the same time, planning quality remains sensitive to context selection, prior experience, and iterative self-correction. Repository-local evidence for the current project suggests that three design directions are central to its implementation: retrieval-style memory augmentation, adaptive capability weighting, and reflection-driven prompt adjustment.

The memory side of the project is closer to case-based reasoning (CBR) and retrieval-augmented generation than to standalone prompting. The codebase stores structured cases, retrieves them through a FAISS-backed semantic index when available, and falls back to keyword retrieval otherwise \citep{lewis2020rag,karpukhin2020dpr,johnson2019faiss,reimers2019sbert}. In broader architectural terms, that places it closer to external-memory and memory-network families \citep{graves2014ntm,weston2015memorynetworks,sukhbaatar2015endtoendmemory} and to classical CBR formulations \citep{aamodt1994cbr,kolodner1993cbr}. The reflection side is closer to iterative self-improvement and agent-memory schemes, including Reflexion, Self-Refine, generative-agent memory, prompt optimization, and the recent Memento line, although the current implementation should still be interpreted strictly as a repository-specific planning loop rather than as a reimplementation of any single external framework \citep{shinn2023reflexion,madaan2023selfrefine,guo2024promptoptimizer,zhou2025memento,park2023generativeagents}.

Within this repository, the target problem is not free-form text generation. Instead, the system aims to produce structured multi-domain operation plans that can be simulated, compared under ablation settings, and exported to architecture and interoperability artifacts. Throughout this manuscript, the term ``joint-operation'' is used in a domain-level sense to refer to coordinated actions across land, maritime, air, and cyber/electronic-warfare domains, rather than in the narrower institutional sense of multi-service joint command (e.g., Joint Chiefs of Staff). This usage reflects the seven-capability-dimension representation (fires, mobility, protection, awareness, EW, sustainment, C2) that spans multiple operational domains within the repository's simulation framework. The five generalization scenarios (coastal assault, urban defense, mountain reconnaissance, river crossing, and maritime sustainment) each engage a different subset of these domains and capability dimensions, and the pipeline is designed to handle this diversity through domain-agnostic structured representations rather than service-specific doctrines. The simulation layer scores plans through a five-stage kill-chain abstraction, computes mission-level metrics such as mission success and overall effectiveness, and aggregates repeated stochastic runs into Monte Carlo summaries. At an operational-reading level, the five stages loosely map to target discovery, disruption and suppression, breach of defended nodes, control of key objectives, and sustainment of tempo and support. This mapping is intended only as a theory-alignment aid: the implemented simulator is still a repository-local operational simplification inspired by joint-targeting, F2T2EA, and network-centric command thinking rather than a doctrinal one-to-one reproduction \citep{jointchiefs2013jointtargeting,tirpak2000f2t2ea,alberts1999networkcentric,law2015simulation}. The export layer produces DoDAF and C2SIM outputs aligned with repository-local structured plan objects and broader simulation-interoperability practice \citep{dodaf2010,ieee2010hla,siso2020c2sim,blais2021automated}.

From a research-positioning perspective, the present study should be read as an application-oriented integration paper rather than as a claim that any one of the three ingredients is individually new. The non-trivial question addressed here is whether these ingredients can be combined into a traceable planning loop in which memory supports context selection, Hope reshapes scenario-conditioned capability emphasis, and reflection repairs iteration-level planning failures without breaking the structured simulation and export chain. The contribution is therefore the implemented coupling and its measured behavior inside a repository-local simulation setting, not a claim of a new standalone planning theory.

This draft intentionally stays close to what can be verified from files in the current project. It does not claim formal guarantees, external benchmark coverage, or a complete literature survey. Instead, the paper makes three application-oriented contributions grounded in code and experiment outputs:

\begin{enumerate}
\item It documents the repository's planning pipeline as an integrated workflow that combines case retrieval, adaptive weighting, bottleneck detection, reflection, structured simulation evaluation, and standards-oriented exporters.
\item It reports a single-scenario five-way ablation (Pure LLM, +Memento, +Hope, +Reflection, Full) using the repository's canonical ablation suite with optimized HOPE parameters, plus a multi-seed replication across five seeds that confirms statistical significance.
\item It reports a five-scenario cross-scenario transfer study using the repository's canonical scene-out evaluation suite and uses that evidence to characterize where gains are stable, where they remain scenario-dependent, and where the current pipeline still appears fragile.
\item It provides a hyperparameter sensitivity analysis for the HOPE controller and a case-bank size ablation that together inform practical deployment choices.
\end{enumerate}

The rest of the paper describes the implemented pipeline, the selected evidence scope, the main quantitative findings, and the current threats to validity that still limit stronger claims.
